% GENERATED FILE. Edit canonical lessons, then run npm run generate:books.
% canonical-source-hash: fnv1a64:ffc42a7657852cf0
% canonical-lessons: KA-C87-really, KA-C87-nothing, KA-C87-ohno, KA-C87-isthatso, KA-C87-thesame

\chapter{Quick Replies}
\label{ch:ka-quick-replies}
\hlchaptermodality{87}

\begin{chapteropening}
\textbf{By the end of this chapter:} \emph{I can say really?, nothing, oh no!, is that so? and the same one.}

Complete the last lesson of chapter 87: I can say really?, nothing, oh no!, is that so? and the same one.
\end{chapteropening}

\section[\texorpdfstring{\kn{ನಿಜವಾಗಲೂ} (nijavāgalū)}{ನಿಜವಾಗಲೂ (nijavāgalū)}]{\kn{ನಿಜವಾಗಲೂ} (nijavāgalū) — really?}

\label{lesson:KA-C87-really}



\begin{quote}
Before the new one: say the Kannada for I didn't hear, then the Kannada for I didn't understand.
\end{quote}

\subsection*{You'll want to know: \kn{ನಿಜವಾಗಲೂ}}
\textbf{\kn{ನಿಜವಾಗಲೂ}?} (\emph{nijavāgalū?}) — \textquotedblleft{}really?\textquotedblright{}, built on \textbf{\kn{ನಿಜ}} (\emph{nija}), \textquotedblleft{}true\textquotedblright{}.

Said with rising voice it is surprise; said flat, it is \textquotedblleft{}truly\textquotedblright{}.

\begin{cousinweb}
\textbf{\kn{ನಿಜವಾಗಲೂ}} (\emph{nijavāgalū}) is \emph{nija}, \textquotedblleft{}true\textquotedblright{}, from Sanskrit \emph{nija}, \textquotedblleft{}innate, one's own\textquotedblright{} (CDIAL 7181), with forms of the native \emph{āgu}, \textquotedblleft{}to become\textquotedblright{} (DEDR 333), and \emph{-ū}, \textquotedblleft{}even\textquotedblright{}: \textquotedblleft{}even truly?\textquotedblright{} DEDR 333 compares Sanskrit \emph{ām}.
\end{cousinweb}

\begin{grammarlens}[title={What you've built}]
True, turned into a question.
\end{grammarlens}

\subsection*{Your turn}
\begin{itemize}
\raggedright
  \item \emph{Say it:} \emph{nijavāgalū}
  \item \emph{Say it:} \emph{nijavāgalū}, once more
  \item \emph{Say it:} say \emph{nijavāgalū?}
  \item \emph{Recall:} say the Kannada for I didn't hear, then the Kannada for I didn't understand, then say \emph{nijavāgalū} again
\end{itemize}

\begin{tcolorbox}[breakable,colback=teal!4,colframe=teal!35!black,title={Before you move on}]
What does \kn{ನಿಜವಾಗಲೂ} mean? (\textquotedblleft{}Really?\textquotedblright{}.) Say it once more.
\end{tcolorbox}
\section[\texorpdfstring{\kn{ಏನೂ} \kn{ಇಲ್ಲ} (ēnū illa)}{ಏನೂ ಇಲ್ಲ (ēnū illa)}]{\kn{ಏನೂ} \kn{ಇಲ್ಲ} (ēnū illa) — nothing}

\label{lesson:KA-C87-nothing}



\begin{quote}
Before the new one: say the Kannada for I didn't understand, then the Kannada for really.
\end{quote}

\subsection*{You'll want to know: \kn{ಏನೂ} \kn{ಇಲ್ಲ}}
\textbf{\kn{ಏನೂ} \kn{ಇಲ್ಲ}} (\emph{ēnū illa}) — \textquotedblleft{}nothing\textquotedblright{}, literally \textquotedblleft{}anything, there isn't\textquotedblright{}. \textbf{\kn{ಏನು}} (\emph{ēnu}), \textquotedblleft{}what\textquotedblright{}, takes \textbf{‑\kn{ಊ}} to become \textquotedblleft{}anything\textquotedblright{}.

\begin{cousinweb}
\textbf{\kn{ಏನೂ} \kn{ಇಲ್ಲ}} (\emph{ēnū illa}) is native Dravidian: \emph{ēnū}, \textquotedblleft{}even what\textquotedblright{}, from \emph{ēnu}, \textquotedblleft{}what\textquotedblright{} (DEDR 5151), and \emph{illa}, \textquotedblleft{}is not\textquotedblright{} (DEDR 2559). \textquotedblleft{}Not even anything\textquotedblright{} is nothing.
\end{cousinweb}

\begin{grammarlens}[title={What you've built}]
What, and not, make nothing.
\end{grammarlens}

\subsection*{Your turn}
\begin{itemize}
\raggedright
  \item \emph{Say it:} \emph{ēnū illa}
  \item \emph{Say it:} \emph{ēnū illa}, once more
  \item \emph{Say it:} answer \emph{ēnū illa}
  \item \emph{Recall:} say the Kannada for I didn't understand, then the Kannada for really, then say \emph{ēnū illa} again
\end{itemize}

\begin{tcolorbox}[breakable,colback=teal!4,colframe=teal!35!black,title={Before you move on}]
What does \kn{ಏನೂ} \kn{ಇಲ್ಲ} mean? (\textquotedblleft{}Nothing\textquotedblright{}.) Say it once more.
\end{tcolorbox}
\section[\texorpdfstring{\kn{ಅಯ್ಯೋ} (ayyō)}{ಅಯ್ಯೋ (ayyō)}]{\kn{ಅಯ್ಯೋ} (ayyō) — oh no!}

\label{lesson:KA-C87-ohno}



\begin{quote}
Before the new one: say the Kannada for really, then the Kannada for nothing.
\end{quote}

\subsection*{You'll want to know: \kn{ಅಯ್ಯೋ}}
\textbf{\kn{ಅಯ್ಯೋ}} (\emph{ayyō}) — \textquotedblleft{}oh no!\textquotedblright{}, \textquotedblleft{}oh dear!\textquotedblright{}. You hear it everywhere in the south when something goes wrong.

\begin{cousinweb}
\textbf{\kn{ಅಯ್ಯೋ}} (\emph{ayyō}) is native Dravidian (DEDR 196). DEDR lists \emph{ayyō}, an interjection of grief, under \emph{ayya}, \textquotedblleft{}father, master\textquotedblright{}. DEDR 196 compares, with a query, Sanskrit \emph{ārya} and \emph{aye}.
\end{cousinweb}

\begin{grammarlens}[title={What you've built}]
The sound of a small disaster.
\end{grammarlens}

\subsection*{Your turn}
\begin{itemize}
\raggedright
  \item \emph{Say it:} \emph{ayyō}
  \item \emph{Say it:} \emph{ayyō}, once more
  \item \emph{Say it:} say \emph{ayyō!}
  \item \emph{Recall:} say the Kannada for really, then the Kannada for nothing, then say \emph{ayyō} again
\end{itemize}

\begin{tcolorbox}[breakable,colback=teal!4,colframe=teal!35!black,title={Before you move on}]
What does \kn{ಅಯ್ಯೋ} mean? (\textquotedblleft{}Oh no!\textquotedblright{}.) Say it once more.
\end{tcolorbox}
\section[\texorpdfstring{\kn{ಹೌದಾ} (haudā)}{ಹೌದಾ (haudā)}]{\kn{ಹೌದಾ} (haudā) — is that so?}

\label{lesson:KA-C87-isthatso}



\begin{quote}
Before the new one: say the Kannada for nothing, then the Kannada for oh no.
\end{quote}

\subsection*{You'll want to know: \kn{ಹೌದಾ}}
\textbf{\kn{ಹೌದಾ}?} (\emph{haudā?}) — \textquotedblleft{}is that so?\textquotedblright{}, \textquotedblleft{}really?\textquotedblright{}. It is \textbf{\kn{ಹೌದು}} (\emph{haudu}), \textquotedblleft{}yes\textquotedblright{}, with the question ending \textbf{‑\kn{ಆ}}.

\begin{cousinweb}
\textbf{\kn{ಹೌದಾ}} (\emph{haudā}) is \emph{haudu}, \textquotedblleft{}yes, it is so\textquotedblright{}, with \emph{-ā}, the ending that makes a question. \emph{Haudu} is usually explained as a form of \emph{āgu}, \textquotedblleft{}to become\textquotedblright{} (DEDR 333). DEDR 333 compares Sanskrit \emph{ām}.
\end{cousinweb}

\begin{grammarlens}[title={What you've built}]
Yes, turned into a question.
\end{grammarlens}

\subsection*{Your turn}
\begin{itemize}
\raggedright
  \item \emph{Say it:} \emph{haudā}
  \item \emph{Say it:} \emph{haudā}, once more
  \item \emph{Say it:} say \emph{haudā?}
  \item \emph{Recall:} say the Kannada for nothing, then the Kannada for oh no, then say \emph{haudā} again
\end{itemize}

\begin{tcolorbox}[breakable,colback=teal!4,colframe=teal!35!black,title={Before you move on}]
What does \kn{ಹೌದಾ} mean? (\textquotedblleft{}Is that so?\textquotedblright{}.) Say it once more.
\end{tcolorbox}
\section[\texorpdfstring{\kn{ಅದೇ} (adē)}{ಅದೇ (adē)}]{\kn{ಅದೇ} (adē) — the same one, that very one}

\label{lesson:KA-C87-thesame}



\begin{quote}
Before the new one: say the Kannada for oh no, then the Kannada for is that so.
\end{quote}

\subsection*{You'll want to know: \kn{ಅದೇ}}
\textbf{\kn{ಅದೇ}} (\emph{adē}) — \textquotedblleft{}that very one, the same\textquotedblright{}. It is \textbf{\kn{ಅದು}} (\emph{adu}), \textquotedblleft{}that\textquotedblright{}, with the emphatic \textbf{‑\kn{ಏ}}.

\textbf{\kn{ಅದೇ}!} — \textquotedblleft{}that's it! exactly that one!\textquotedblright{}

\begin{cousinweb}
\textbf{\kn{ಅದೇ}} (\emph{adē}) is native Dravidian: \emph{adu}, \textquotedblleft{}that\textquotedblright{} (DEDR 1), with \emph{-ē}, the ending that stresses: \textquotedblleft{}that very one\textquotedblright{}.
\end{cousinweb}

\begin{grammarlens}[title={What you've built}]
That, made emphatic.
\end{grammarlens}

\subsection*{Your turn}
\begin{itemize}
\raggedright
  \item \emph{Say it:} \emph{adē}
  \item \emph{Say it:} \emph{adē}, once more
  \item \emph{Say it:} say \emph{adē!}
  \item \emph{Recall:} say the Kannada for oh no, then the Kannada for is that so, then say \emph{adē} again
\end{itemize}

\begin{tcolorbox}[breakable,colback=teal!4,colframe=teal!35!black,title={Before you move on}]
What does \kn{ಅದೇ} mean? (\textquotedblleft{}The same one, that very one\textquotedblright{}.) Say it once more.
\end{tcolorbox}
